\section{Explaining the lack of prominence}

The preceding sections indicate that citizens are not uniformly opposed to global redistributive and climate policies. % AF: c'est bcp trop conservateur comme restitution: majorities generally accept well-designed internationally redistributive policies semble une représentation plus fidèle des résultats
On the contrary, several such policies attract conditional but frequently majority support, including when their costs and institutional design are made explicit. This creates a %central AF: delete
puzzle: why do apparently acceptable policies remain marginal in electoral competition and public debate \citep{fabre_majority_2025,fabre_public_2026}? The finding is not confined to direct survey questions. Support remains visible, although its magnitude varies, in list experiments, petition decisions, conjoint evaluations, and policy-prioritization tasks \citep{fabre_majority_2025}.

Addressing this puzzle requires distinguishing public support from political prominence and issue salience. Political prominence concerns an issue's visibility in party platforms, electoral campaigns, and public debate, whereas issue salience captures the importance citizens assign to that issue and the weight it carries in their political decisions. A policy may therefore attract majority support when evaluated in isolation without becoming sufficiently important to influence candidate evaluations or vote choice. Research on attitude importance shows that policy preferences exert greater influence on electoral behavior when citizens consider the underlying issue personally important \citep{krosnick_role_1988,dennison_review_2019}.

Beyond limited electoral salience, % AF: "limited electoral salience" is not substantiated
the gap between first-order preferences and second-order beliefs provides another potential explanation. As shown in the preceding section, citizens' reported willingness to contribute to climate action substantially exceeds their estimates of others' willingness to do so \citep{andre_globally_2024}. These misperceptions can create a coordination problem: individuals may refrain from expressing or acting on a favorable preference when they expect little participation from others. Consistent with this mechanism, \citet{mildenberger_beliefs_2017} find that members of the public and political elites systematically underestimate the prevalence of pro-climate positions and that experimentally correcting these second-order beliefs increases support for climate-policy action.

A similar mechanism may constrain political supply. \citet{fabre_public_2026} suggests that the limited supply of global redistributive proposals may reflect pluralistic ignorance among policymakers and activists, who may underestimate universalistic attitudes, public support, or the electoral returns to endorsing such policies. More general evidence shows that political representatives frequently misperceive constituency opinion and often infer that voters are more conservative than they actually are \citep{broockman_bias_2018,pilet_politicians_2023}. Although these studies do not directly examine global redistribution, they provide evidence for the plausibility of this supply-side mechanism.

% AF: ce qui serait intéressant ici c'est de comparer les expériences de donation de Nair (Respondents in the control group elected to keep 61.7% of their bonus, while donating 28.0% to a domestic charity and donating 10.1% to an international charity.) et budget allocation task de NHB et Fabre 26, ainsi que les expérience donation de Fehr et NHB (et d'autres s'il y en a d'autres). J'ai pas bcp cherché mais pas sûr que Fehr donnent leur stats descriptives (pour comparer les distributions -ou au moins les moyennes- de dons quand le destinataire est un compatriote ou un étranger): leur demander si elles ne sont pas dans le papier. 
Information alone does not necessarily transform favorable attitudes into political demand. In the United States, correcting misperceptions of relative global affluence increases support for foreign economic assistance and raises real-stakes donations to international charities, but produces only a modest increase in willingness to petition elected representatives \citep{nair_misperceptions_2018}. Respondents also continue to donate substantially more to domestic than to international charities. Similarly, although Germans systematically underestimate their position in the global income distribution, correcting this misperception affects neither support for global redistribution nor donations to the global poor \citep{fehr_your_2019}. These findings indicate that national reference groups may remain more politically influential than comparisons with the global income distribution.

Material circumstances can reinforce this national framing. Economically vulnerable left-leaning citizens may consider assistance to poorer countries morally justified while remaining reluctant to increase public expenditure when foreign aid is perceived as competing with domestic priorities \citep{cho_ideology_2024}. % AF: ça veut dire quoi "may"? on veut des résultas, des chiffres
% Misperceptions of existing aid cannot fully explain this reluctance. Citizens in donor countries substantially overestimate current aid expenditure but, on average, still prefer a level exceeding both actual and perceived aid \citep{kaufmann_foreign_2019}. % AF: c'est pour la section 2 ça

The limited prominence of international redistribution therefore cannot be reduced to widespread public opposition. % AF: c'est pas le sujet ici, cette section c'est sur pourquoi il y a une faible visibilité de la redistribution mondiale malgré son soutien majoritaire
It instead reflects the interaction of low electoral salience, pessimistic beliefs about others' support, nationally bounded reference groups, material constraints, and a political supply that may itself respond to biased perceptions and unequal political influence.

